\section{Q2: CNN vs.\ MLP, spatial structure, and local vs.\ global information}
\label{sec:q2}

\paragraph{Reuse.} Q2 compares \textbf{MLP-M} and the \textbf{CNN} from Q1(a) (Tab.~\ref{tab:models}). Their clean-image runs
and fixed-permutation runs were first trained in Q1(a) and are loaded from its run cache with the same seeds, LR grid and
checkpoints, so parts (a) and (b) reuse those results. Part (c) is new.

\subsection{(a) A fair comparison}
\paragraph{Fairness criterion (fixed before any result).} Both models (1) have the same number of trainable parameters,
within 0.2\%; (2) see identical data and normalisation, with no augmentation; (3) get an identical budget and recipe, \ie
30 epochs at batch 128, AdamW, cosine LR, weight decay 0.05, BatchNorm and ReLU; (4) get the same LR grid, with the LR and
early-stopping epoch chosen on validation only; and (5) are trained with 3 paired seeds. The primary outcome is mean$\pm$sd
test accuracy, evaluated once. Compute per image is declared unmatched ($\approx$118$\times$ more MACs for the CNN) and is
controlled by part (b).

\begin{table}[t]
\centering\footnotesize\setlength{\tabcolsep}{3pt}
\caption{Q2: the compared models (both introduced in Q1(a)).}
\label{tab:q2models}
\begin{tabular}{@{}lrrll@{}}
\toprule
Model & Params & MACs & LR (clean) & Test acc.\ (clean, \%) \\
\midrule
MLP-M & \res{q1.MLP-M.params} & \res{q1.MLP-M.macs} & \res{q2.lr.MLP-M.clean} & \res{q2.MLP-M.clean.test} \\
CNN   & \res{q1.CNN.params}   & \res{q1.CNN.macs}   & \res{q2.lr.CNN.clean}   & \res{q2.CNN.clean.test} \\
\bottomrule
\end{tabular}
\end{table}

\paragraph{Proposed reason.} The training curves (Fig.~\ref{fig:q1curves}) show the MLP fitting its training set while
generalising worse. Its first-layer units are whole-image, position-specific templates, while the CNN's are small local
filters reused everywhere (Fig.~\ref{fig:q2weights}). We therefore propose: \emph{locality and weight sharing let the CNN
exploit the correlations between neighbouring pixels with a small set of reusable local detectors, which the MLP must
relearn separately at every position}.

\begin{figure}[t]
\centering
\resfig{fig_q2_2_first_layer_weights.png}{\linewidth}{1.0in}
\caption{First-layer weights. Left: the 12 largest-norm MLP-M units (each spans the whole image). Right: all 32 CNN
$3{\times}3$ filters.}
\label{fig:q2weights}
\end{figure}

\subsection{(b) Same pixel values, spatial structure destroyed}
We apply the fixed global pixel permutation of Q1(a) to the training, validation and test images; every image keeps
exactly the same multiset of RGB values (Tab.~\ref{tab:q2tr}), and both models are retrained. \textbf{Expected beforehand:}
MLP-M unchanged within seed noise ($|\Delta|\leq$ 2 seed-sd, at least 1\,pp), because the permutation is absorbed by its
first layer; CNN drops to about the MLP's level or below (the gap shrinks by more than half), because its $3{\times}3$ filters
now see unrelated pixels. Parameters, FLOPs and pixel values are unchanged, so if the CNN kept its advantage, the reason
proposed in (a) would be wrong.

\subsection{(c) Local structure or global arrangement?}
A full permutation destroys all spatial structure at once, so it cannot separate two explanations of \emph{what} the CNN
uses. \textbf{E-local:} local patterns (edges, textures, colour blobs within a few pixels) pooled over the image, \ie a
bag of local features~\cite{brendel2019approximating}; its global average pooling makes this plausible. \textbf{E-global:} the
arrangement of parts across the image, \ie object shape. \textbf{Experiment:} \emph{per-image random tile shuffling} with
tile size $t\in\{16,8,4\}$. Each image is cut into $(32/t)^2$ tiles, which are reassembled in an order drawn separately
for each image (seeded; fixed once, not re-drawn per epoch). Pixels inside a tile keep their neighbours, while the
global arrangement becomes random and carries no information. The test images go through the same procedure with the
same $t$. Tab.~\ref{tab:q2tr} quantifies how much structure survives. Both models are retrained under each condition with
the same treatment as in (a), including LR selection on validation.

\paragraph{Predictions and decision rule (pre-registered).} Let the CNN's \emph{retention}
$R(t)=\frac{\mathrm{acc}(t)-\mathrm{acc}(\text{perm})}{\mathrm{acc}(\text{clean})-\mathrm{acc}(\text{perm})}$ be the share of
its spatial benefit that survives. E-local predicts $R(8)\geq0.5$, a CNN--MLP gap that stays large, and confident correct
predictions. E-global predicts $R(8)<0.5$, already low at $t{=}16$. The reason proposed in (a) also predicts that the
\emph{MLP} loses accuracy under tile shuffling (unlike under the fixed permutation), because its templates need content at
consistent positions. Decision: E-local if $R(8)\geq0.5$ on test, otherwise E-global. A secondary check applies the
same normalisation to the CNN's mean confidence on its correct predictions ($R_c$).

\begin{table}[t]
\centering\footnotesize\setlength{\tabcolsep}{2.6pt}
\caption{Transformations (measured on 1,000 training images). ``Same values'': identical multiset of RGB pixels per image.
$p_d$: fraction of pixel pairs at distance $d$ that remain at distance $d$ in the same direction.}
\label{tab:q2tr}
\begin{tabular}{@{}lcrrrrr@{}}
\toprule
Condition & Same values & $p_1$ & $p_2$ & $p_4$ & $p_8$ & $p_{16}$ \\
\midrule
clean  & \res{q2.tr.clean.same}  & \res{q2.tr.clean.d1}  & \res{q2.tr.clean.d2}  & \res{q2.tr.clean.d4}  & \res{q2.tr.clean.d8}  & \res{q2.tr.clean.d16} \\
perm (b) & \res{q2.tr.perm.same} & \res{q2.tr.perm.d1}   & \res{q2.tr.perm.d2}   & \res{q2.tr.perm.d4}   & \res{q2.tr.perm.d8}   & \res{q2.tr.perm.d16} \\
tile 16 & \res{q2.tr.tile16.same} & \res{q2.tr.tile16.d1} & \res{q2.tr.tile16.d2} & \res{q2.tr.tile16.d4} & \res{q2.tr.tile16.d8} & \res{q2.tr.tile16.d16} \\
tile 8  & \res{q2.tr.tile8.same}  & \res{q2.tr.tile8.d1}  & \res{q2.tr.tile8.d2}  & \res{q2.tr.tile8.d4}  & \res{q2.tr.tile8.d8}  & \res{q2.tr.tile8.d16} \\
tile 4  & \res{q2.tr.tile4.same}  & \res{q2.tr.tile4.d1}  & \res{q2.tr.tile4.d2}  & \res{q2.tr.tile4.d4}  & \res{q2.tr.tile4.d8}  & \res{q2.tr.tile4.d16} \\
\bottomrule
\end{tabular}
\end{table}

\begin{figure}[t]
\centering
\resfig{fig_q2_0_conditions.png}{0.78\linewidth}{2.3in}
\caption{The same training images under each condition.}
\label{fig:q2cond}
\end{figure}

\begin{table}[t]
\centering\footnotesize\setlength{\tabcolsep}{3pt}
\caption{\textbf{Q2 test accuracy} (\%, mean$\pm$sd over 3 seeds) for MLP-M (\res{q1.MLP-M.params} params) and the CNN
(\res{q1.CNN.params} params). Gap: CNN $-$ MLP (pp). $R$: CNN retention of its spatial benefit.}
\label{tab:q2test}
\begin{tabular}{@{}llrrrr@{}}
\toprule
Part & Condition & MLP-M & CNN & Gap & $R$ \\
\midrule
(a) & clean & \res{q2.MLP-M.clean.test} & \res{q2.CNN.clean.test} & \res{q2.gap.clean} & 1 \\
(b) & fixed permutation & \res{q2.MLP-M.perm.test} & \res{q2.CNN.perm.test} & \res{q2.gap.perm} & 0 \\
(c) & tiles, $t{=}16$ & \res{q2.MLP-M.tile16.test} & \res{q2.CNN.tile16.test} & \res{q2.gap.tile16} & \res{q2.R.tile16} \\
(c) & tiles, $t{=}8$  & \res{q2.MLP-M.tile8.test}  & \res{q2.CNN.tile8.test}  & \res{q2.gap.tile8}  & \res{q2.R.tile8} \\
(c) & tiles, $t{=}4$  & \res{q2.MLP-M.tile4.test}  & \res{q2.CNN.tile4.test}  & \res{q2.gap.tile4}  & \res{q2.R.tile4} \\
\bottomrule
\end{tabular}
\end{table}

\begin{figure}[t]
\centering
\resfig{fig_q2_5_tile_curve.png}{0.92\linewidth}{1.9in}
\caption{Test accuracy as the global arrangement is removed (dashed: fixed permutation of part (b)).}
\label{fig:q2tile}
\end{figure}

\begin{table}[t]
\centering\footnotesize\setlength{\tabcolsep}{2.4pt}
\caption{\textbf{Q2 test confidence and calibration} (max-softmax; means over seeds). Conf\,c/w: mean confidence of
correct/wrong predictions. ECE: expected calibration error (15 bins)~\cite{guo2017calibration}. AUROC: confidence as a
detector of correct predictions.}
\label{tab:q2conf}
\begin{tabular}{@{}llrrrrr@{}}
\toprule
Model & Condition & Conf\,c & Conf\,w & ECE & AUROC & NLL \\
\midrule
MLP-M & clean & \res{q2.MLP-M.clean.confc} & \res{q2.MLP-M.clean.confw} & \res{q2.MLP-M.clean.ece} & \res{q2.MLP-M.clean.auroc} & \res{q2.MLP-M.clean.nll} \\
MLP-M & tile 8 & \res{q2.MLP-M.tile8.confc} & \res{q2.MLP-M.tile8.confw} & \res{q2.MLP-M.tile8.ece} & \res{q2.MLP-M.tile8.auroc} & \res{q2.MLP-M.tile8.nll} \\
MLP-M & perm & \res{q2.MLP-M.perm.confc} & \res{q2.MLP-M.perm.confw} & \res{q2.MLP-M.perm.ece} & \res{q2.MLP-M.perm.auroc} & \res{q2.MLP-M.perm.nll} \\
CNN & clean & \res{q2.CNN.clean.confc} & \res{q2.CNN.clean.confw} & \res{q2.CNN.clean.ece} & \res{q2.CNN.clean.auroc} & \res{q2.CNN.clean.nll} \\
CNN & tile 16 & \res{q2.CNN.tile16.confc} & \res{q2.CNN.tile16.confw} & \res{q2.CNN.tile16.ece} & \res{q2.CNN.tile16.auroc} & \res{q2.CNN.tile16.nll} \\
CNN & tile 8 & \res{q2.CNN.tile8.confc} & \res{q2.CNN.tile8.confw} & \res{q2.CNN.tile8.ece} & \res{q2.CNN.tile8.auroc} & \res{q2.CNN.tile8.nll} \\
CNN & tile 4 & \res{q2.CNN.tile4.confc} & \res{q2.CNN.tile4.confw} & \res{q2.CNN.tile4.ece} & \res{q2.CNN.tile4.auroc} & \res{q2.CNN.tile4.nll} \\
CNN & perm & \res{q2.CNN.perm.confc} & \res{q2.CNN.perm.confw} & \res{q2.CNN.perm.ece} & \res{q2.CNN.perm.auroc} & \res{q2.CNN.perm.nll} \\
\bottomrule
\end{tabular}
\end{table}

\begin{figure*}[t]
\centering
\resfig{fig_q2_6_reliability.png}{0.95\textwidth}{1.5in}\\[2pt]
\resfig{fig_q2_7_confidence_hist.png}{0.66\textwidth}{2.2in}
\caption{Q2 test-set confidence analysis. Top: reliability diagrams per condition. Bottom: confidence of correct vs.\
wrong predictions (clean, $t{=}8$, permuted). Per-class accuracy and CNN confusion matrices are in
\texttt{results/q2\_outputs/fig\_q2\_8*}, \texttt{fig\_q2\_9*}.}
\label{fig:q2conf}
\end{figure*}

\paragraph{Results and conclusion.} The CNN's confidence retention is $R_c(8)=$ \res{q2.Rc.tile8}. The paragraphs below
were generated from Tabs.~\ref{tab:q2test}--\ref{tab:q2conf} by the pre-registered rules (Q2 notebook \S6).

\noindent\textbf{(a)} \restext{q2a}

\noindent\textbf{(b)} \restext{q2b}

\noindent\textbf{(c)} \restext{q2c}

\noindent\textbf{Linking the three experiments.} \restext{q2link}
